% BHU PYQ -- Manually transcribed & error-corrected
% Paper: BCF-125 : Fundamentals of Banking
% Exam: B.Com. (Hons.) FMM Semester II Examination, 2021-22

\documentclass[11pt,a4paper]{article}
\usepackage[a4paper,top=2.5cm,bottom=2.5cm,left=2.8cm,right=2.8cm,headheight=14pt]{geometry}
\usepackage{amsmath,amssymb}
\usepackage[T1]{fontenc}
\usepackage[utf8]{inputenc}
\usepackage{lmodern,microtype,fancyhdr,enumitem,hyperref,lastpage,parskip}

\hypersetup{
    colorlinks=true,
    urlcolor=black,
    linkcolor=black,
    pdftitle={BCF-125: Fundamentals of Banking -- BHU B.Com. (Hons.) FMM Sem II 2021-22}
}

\pagestyle{fancy}
\fancyhf{}
\renewcommand{\headrulewidth}{0.4pt}
\renewcommand{\footrulewidth}{0.4pt}
\fancyhead[L]{\small   \textit{Banaras Hindu University}}
\fancyhead[R]{\small   \textit{BCF-125: Fundamentals of Banking}}
\fancyfoot[C]{\small Page  \thepage\ of \pageref{LastPage}}


\newlist{parts}{enumerate}{2}
\setlist[parts,1]{label=(\alph*),leftmargin=*,itemsep=5pt,topsep=3pt}
\setlist[parts,2]{label=(\roman*),leftmargin=*,itemsep=3pt,topsep=2pt}

\newcommand{\pts}[1]{\hfill   \textit{[#1]}}


\begin{document}


\begin{center}
  {\large\bfseries BANARAS HINDU UNIVERSITY}\\[2pt]
  {\normalsize B.Com. (Hons.) FMM --- Semester II Examination, 2021-22}\\[6pt]
  \rule{0.8\linewidth}{0.5pt}\\[6pt]
  {\large\bfseries Paper No. BCF-125: Fundamentals of Banking}\\[4pt]
  {\normalsize Time: 3 Hours\hfill Maximum Marks: 70}
\end{center}

\rule{\linewidth}{0.4pt}

\smallskip



\noindent   \textit{Instructions: Answer all questions. All questions carry equal marks (14 marks each).}

\bigskip




\noindent   \textbf{Question 1.}
What do you understand by Commercial Bank? Explain the concept of customer-bank relationship.\pts{14}

\centerline{   \textbf{OR}}

Give a brief note on the role of banking in the Indian Economy.\pts{14}

\medskip\rule{\linewidth}{0.2pt}




\noindent   \textbf{Question 2.}
Describe the different types of banks and explain their functions.\pts{14}

\centerline{   \textbf{OR}}

What is a Central Bank? Differentiate between a Commercial Bank and a Central Bank.\pts{14}

\medskip\rule{\linewidth}{0.2pt}




\noindent   \textbf{Question 3.}
Explain the process of deposit mobilization of a Commercial Bank.\pts{14}

\centerline{   \textbf{OR}}

Explain the role of the Reserve Bank of India (RBI) in the credit control process.\pts{14}

\medskip\rule{\linewidth}{0.2pt}




\noindent   \textbf{Question 4.}
What is priority sector lending? Explain the targets and categories under priority sector lending.\pts{14}

\centerline{   \textbf{OR}}

Explain the different kinds of charge over securities.\pts{14}

\medskip\rule{\linewidth}{0.2pt}




\noindent   \textbf{Question 5.}
Discuss the concept of Development Banks and describe the major problems faced by development banks in the post-liberalisation era.\pts{14}

\centerline{   \textbf{OR}}

Give brief notes on any two of the following:\pts{14}

\begin{parts}
  \item NABARD
  \item ICICI
  \item SIDBI
\end{parts}

\vfill
\rule{\linewidth}{0.4pt}

\begin{center}\small
\href{https://bhu-exam.vercel.app/}{Click here to visit our website}\\[4pt]
\href{https://me-aryan.vercel.app/}{Click here to visit our contributor}\\[4pt]
\href{https://quashan-me.vercel.app/}{Click here to visit our contributor}
\end{center}

\end{document}
